\chapter{FUTURE ENHANCEMENT}

\section{Limitations}

Despite achieving the intended objectives, several limitations were identified in the proposed system.

First, the model size was restricted to the 7B parameter configuration due to hardware constraints. Although QLoRA significantly reduces memory requirements, larger models could potentially provide improved reasoning capabilities and higher-quality code generation.

Second, the dataset size and diversity were limited. While preprocessing and instruction formatting improved data quality, the dataset lacks broader domain coverage, which may affect the model’s generalization to unseen or complex programming tasks.

Third, evaluation relied heavily on automated metrics such as BLEU and CodeBLEU. Although useful for measuring syntactic similarity, these metrics do not fully capture functional correctness or logical soundness. Execution-based evaluation was performed, but only on a limited subset of test cases.

Additionally, a mismatch exists between training and evaluation data formats. The HumanEval benchmark provides function signatures in prompts, whereas the training datasets (e.g., Alpaca and Flytech) do not consistently include such structured inputs. This discrepancy likely impacted functional performance scores.

Finally, hyperparameter optimization was conducted under constrained computational budgets, limiting the exploration of the full search space. Long-context reasoning and multi-file code generation capabilities were also not extensively explored.

\section{Future Work}

Several directions can be pursued to further enhance the system and extend its capabilities.

\subsection{Scaling Model Capacity}

Future work may explore fine-tuning larger models (e.g., 13B or 70B variants) or adopting mixture-of-experts architectures to improve reasoning depth and contextual understanding. Investigating scaling laws could help determine optimal trade-offs between model size, dataset size, and computational resources.

\subsection{Dataset Expansion and Curriculum Learning}

Expanding the dataset to include diverse programming scenarios can significantly improve robustness. This includes:
\begin{itemize}
    \item Multi-language programming tasks (C++, Java, Go),
    \item Repository-level and real-world code,
    \item Complex algorithmic and competitive programming problems,
    \item Debugging and refactoring tasks.
\end{itemize}

Incorporating curriculum learning—progressing from simple to complex tasks—may further improve training stability and model performance.

\subsection{Advanced Evaluation Strategies}

Future evaluation frameworks should incorporate:
\begin{itemize}
    \item Large-scale execution-based testing,
    \item Static code analysis tools,
    \item Security and vulnerability detection metrics,
    \item Human expert evaluation.
\end{itemize}

These approaches would provide a more comprehensive understanding of model reliability and practical usability.

\subsection{Retrieval-Augmented Code Generation}

Integrating retrieval-based mechanisms can enable the model to access external knowledge sources such as documentation, APIs, and code repositories. This hybrid approach can improve factual accuracy and real-world applicability.

\subsection{Deployment and System Integration}

Future enhancements may focus on real-world deployment, including:
\begin{itemize}
    \item Integration as an IDE-based coding assistant,
    \item Optimization for real-time inference,
    \item Quantization-aware serving techniques,
    \item Integration into CI/CD pipelines for automated validation.
\end{itemize}

Such improvements would make the system more practical for production-level applications.